\chapterwithopening{System Analysis and Requirements}{chap:system_analysis}{
  \chaptercontentsitem{sec:system_analysis_introduction}
  \chaptercontentsitem{sec:system_overview}
  \chaptercontentsitem{sec:system_actors}
  \chaptercontentsitem{sec:functional_requirements}
  \chaptercontentsitem{sec:non_functional_requirements}
  \chaptercontentsitem{sec:use_case_model}
  \chaptercontentsitem{sec:activity_diagrams}
  \chaptercontentsitem{sec:chapter4_summary}
}

\tikzset{
  edu activity/.style={
    font=\footnotesize,
    >=Latex,
    node distance=6mm and 16mm,
    every node/.style={align=center}
  },
  edu startstop/.style={
    ellipse,
    draw,
    thick,
    fill=gray!15,
    minimum width=26mm,
    minimum height=7mm,
    inner sep=2pt
  },
  edu process/.style={
    rectangle,
    rounded corners=2pt,
    draw,
    thick,
    fill=blue!6,
    text width=82mm,
    minimum height=8mm,
    inner sep=3pt
  },
  edu decision/.style={
    diamond,
    draw,
    thick,
    fill=orange!10,
    aspect=2.3,
    text width=28mm,
    inner sep=1pt
  },
  edu arrow/.style={
    ->,
    thick
  },
  edu label/.style={
    font=\footnotesize,
    fill=white,
    inner sep=1pt
  }
}

\newsavebox{\EduActivityBox}

\newcommand{\EduActivityFigure}[3]{%
  \begin{figure}[htbp]
    \centering
    \sbox{\EduActivityBox}{#1}%
    \ifdim\wd\EduActivityBox>0.96\textwidth
      \resizebox{0.96\textwidth}{!}{\usebox{\EduActivityBox}}%
    \else\ifdim\dimexpr\ht\EduActivityBox+\dp\EduActivityBox\relax>0.54\textheight
      \resizebox{!}{0.54\textheight}{\usebox{\EduActivityBox}}%
    \else
      \usebox{\EduActivityBox}%
    \fi\fi
    \caption{#2}
    \label{#3}
  \end{figure}
}

\tikzset{
  ch4 act step/.style={
    edu process,
    text width=72mm,
    minimum height=9mm
  },
  ch4 act step medium/.style={
    edu process,
    text width=62mm,
    minimum height=9mm
  },
  ch4 act step branch/.style={
    edu process,
    text width=48mm,
    minimum height=9mm
  },
  ch4 act decision/.style={
    edu decision,
    text width=30mm,
    aspect=2.0
  },
  ch4 act decision wide/.style={
    edu decision,
    text width=34mm,
    aspect=2.0
  },
  ch4 act merge/.style={
    diamond,
    draw=black!85,
    line width=0.9pt,
    fill=white,
    minimum size=7mm,
    inner sep=0pt,
    outer sep=0pt
  },
  ch4 rebuilt card/.style={
    draw=black!35,
    rounded corners=5pt,
    fill=black!1,
    line width=0.75pt
  },
  ch4 rebuilt step/.style={
    edu process,
    text width=38mm,
    minimum height=8.5mm,
    inner sep=3pt
  },
  ch4 rebuilt step wide/.style={
    edu process,
    text width=52mm,
    minimum height=8.5mm,
    inner sep=3pt
  },
  ch4 rebuilt step narrow/.style={
    edu process,
    text width=31mm,
    minimum height=8.2mm,
    inner sep=2.6pt
  },
  ch4 rebuilt decision/.style={
    edu decision,
    text width=25mm,
    aspect=2.05,
    inner sep=1pt
  },
  ch4 rebuilt fork/.style={
    rectangle,
    draw=black!85,
    fill=black!85,
    minimum height=0.9mm,
    inner sep=0pt,
    outer sep=0pt
  },
  ch4 rebuilt title/.style={
    font=\footnotesize\bfseries,
    anchor=west,
    align=left
  }
}

\newcommand{\EduActivityAuthWorkspaceAccessDiagram}{%
\begin{tikzpicture}[edu activity, x=1mm, y=1mm]
  \node (start) [edu initial] at (0,0) {};
  \node (open) [ch4 rebuilt step wide] at (0,-12) {User opens the EduVerse access page.};
  \node (choice) [ch4 rebuilt decision] at (0,-31) {Choose access path?};

  \draw[ch4 rebuilt card] (-78,-48) rectangle (-24,-126);
  \node[ch4 rebuilt title] at (-70,-55) {Sign up};
  \node (signup) [ch4 rebuilt step narrow] at (-49,-69) {Register a new account.};
  \node (verify) [ch4 rebuilt step narrow] at (-49,-91) {Confirm the verification email.};
  \node (signupreturn) [ch4 rebuilt step narrow] at (-49,-113) {Return to sign in.};

  \draw[ch4 rebuilt card] (-20,-48) rectangle (20,-82);
  \node[ch4 rebuilt title] at (-16,-55) {Sign in};
  \node (signin) [ch4 rebuilt step narrow] at (0,-72) {Sign in with credentials.};

  \draw[ch4 rebuilt card] (24,-48) rectangle (78,-148);
  \node[ch4 rebuilt title] at (29,-55) {Reset access};
  \node (reset) [ch4 rebuilt step narrow] at (49,-69) {Request password reset.};
  \node (resetmail) [ch4 rebuilt step narrow] at (49,-91) {Receive password-reset email.};
  \node (resetlink) [ch4 rebuilt step narrow] at (49,-113) {Open the reset link.};
  \node (newpass) [ch4 rebuilt step narrow] at (49,-135) {Set the new password.};

  \node (valid) [ch4 rebuilt decision] at (0,-162) {Access accepted?};
  \node (error) [ch4 rebuilt step narrow] at (56,-173) {Error message is shown.};
  \node (errorend) [edu final] at (56,-196) {};
  \node (context) [ch4 rebuilt step wide] at (0,-191) {Select active organization or role if needed.};
  \node (workspace) [ch4 rebuilt step wide] at (0,-214) {Dashboard or workspace is displayed.};
  \node (logout) [ch4 rebuilt step wide] at (0,-237) {User works in the selected context and logs out.};
  \node (end) [edu final] at (0,-258) {};

  \path [edu arrow] (start.south) -- (open.north);
  \path [edu arrow] (open) -- (choice);
  \path [edu arrow] (choice.west) -| (signup.north);
  \path [edu arrow] (choice) -- (signin.north);
  \path [edu arrow] (choice.east) -| (reset.north);
  \path [edu arrow] (signup) -- (verify);
  \path [edu arrow] (verify) -- (signupreturn);
  \path [edu arrow] (reset) -- (resetmail);
  \path [edu arrow] (resetmail) -- (resetlink);
  \path [edu arrow] (resetlink) -- (newpass);
  \path [edu arrow] (signupreturn.east) -- ++(8mm,0) |- (signin.west);
  \path [edu arrow] (signin) -- (valid);
  \path [edu arrow] (newpass.west) -- ++(-8mm,0) |- (signin.east);
  \path [edu arrow] (valid.east) -- node[edu label, above] {No} ++(8mm,0) -| (error.north);
  \path [edu arrow] (error) -- (errorend);
  \path [edu arrow] (valid) -- node[edu label] {Yes} (context);
  \path [edu arrow] (context) -- (workspace);
  \path [edu arrow] (workspace) -- (logout);
  \path [edu arrow] (logout) -- (end);
\end{tikzpicture}%
}

\newcommand{\EduActivityOrganizationAccessDiagram}{%
\resizebox*{!}{0.70\textheight}{%
\begin{tikzpicture}[edu activity, x=1mm, y=1mm]
  \node (start) [edu initial] at (0,0) {};
  \node (open) [ch4 rebuilt step wide] at (0,-12) {Administrator opens the people and access workspace.};
  \node (mode) [ch4 rebuilt decision] at (0,-31) {Select access method.};
  \node (fork) [ch4 rebuilt fork, minimum width=108mm] at (0,-47) {};

  \draw[ch4 rebuilt card] (-74,-58) rectangle (-6,-192);
  \node[ch4 rebuilt title] at (-70,-65) {A. Invitation-based access};
  \node (invite) [ch4 rebuilt step] at (-40,-84) {Administrator invites or registers a member.};
  \node (recipient) [ch4 rebuilt step] at (-40,-111) {User opens the invitation.};
  \node (invitesignin) [ch4 rebuilt step] at (-40,-138) {User signs in or signs up if required.};
  \node (inviteaccess) [ch4 rebuilt step] at (-40,-165) {Organization workspace is opened.};
  \node (inviteend) [edu final] at (-40,-185) {};

  \draw[ch4 rebuilt card] (6,-58) rectangle (74,-266);
  \node[ch4 rebuilt title] at (10,-65) {B. Public Join Link access};
  \node (openlink) [ch4 rebuilt step] at (40,-84) {User opens the public join link.};
  \node (linksignin) [ch4 rebuilt step] at (40,-111) {User signs in or signs up if required.};
  \node (approval) [ch4 rebuilt decision] at (40,-136) {Approval required?};
  \node (directaccess) [ch4 rebuilt step narrow] at (23,-164) {Access is granted.};
  \node (directend) [edu final] at (23,-188) {};
  \node (request) [ch4 rebuilt step narrow] at (58,-164) {User submits a join request.};
  \node (review) [ch4 rebuilt step narrow] at (58,-190) {Administrator reviews the request.};
  \node (approve) [ch4 rebuilt decision] at (58,-214) {Approved?};
  \node (reject) [ch4 rebuilt step narrow] at (23,-234) {Rejection or pending status is shown.};
  \node (grant) [ch4 rebuilt step narrow] at (58,-238) {Organization workspace is opened.};
  \node (rejectend) [edu final] at (23,-255) {};
  \node (grantend) [edu final] at (58,-259) {};

  \path [edu arrow] (start.south) -- (open.north);
  \path [edu arrow] (open) -- (mode);
  \path [edu arrow] (mode) -- (fork);
  \path [edu arrow] (fork.west) -| node[edu label, pos=0.25] {Invitation} (invite.north);
  \path [edu arrow] (fork.east) -| node[edu label, pos=0.25] {Public link} (openlink.north);
  \path [edu arrow] (invite) -- (recipient);
  \path [edu arrow] (recipient) -- (invitesignin);
  \path [edu arrow] (invitesignin) -- (inviteaccess);
  \path [edu arrow] (inviteaccess) -- (inviteend);
  \path [edu arrow] (openlink) -- (linksignin);
  \path [edu arrow] (linksignin) -- (approval);
  \path [edu arrow] (approval.west) -| node[edu label, pos=0.25] {No} (directaccess.north);
  \path [edu arrow] (directaccess) -- (directend);
  \path [edu arrow] (approval.east) -| node[edu label, pos=0.25] {Yes} (request.north);
  \path [edu arrow] (request) -- (review);
  \path [edu arrow] (review) -- (approve);
  \path [edu arrow] (approve.west) -| node[edu label, pos=0.25] {No} (reject.north);
  \path [edu arrow] (reject) -- (rejectend);
  \path [edu arrow] (approve) -- node[edu label] {Yes} (grant);
  \path [edu arrow] (grant) -- (grantend);
\end{tikzpicture}%
}%
}

\newcommand{\EduActivityClassLifecycleDiagram}{%
\begin{tikzpicture}[edu activity, node distance=8.5mm and 28mm]
  \node (start) [edu initial] {};
  \node (open) [ch4 act step] {Administrator or permitted Teacher opens the class-management workspace.};
  \node (role) [ch4 act decision wide, below=of open] {Administrator view or Teacher view?};
  \node (adminview) [ch4 act step branch, left=50mm of role] {Administrator manages organization-wide class setup and oversight.};
  \node (teacherview) [ch4 act step branch, right=50mm of role] {Teacher opens the class controls allowed by organization permissions.};
  \node (details) [ch4 act step, below=16mm of role] {Create or edit the class and define the title, code, term, stage, and other details.};
  \node (members) [ch4 act step, below=of details] {Assign or change the Teacher and manage class membership.};
  \node (features) [ch4 act step, below=of members] {Configure class tools, visibility rules, and Results visibility.};
  \node (save) [ch4 act step medium, below=of features] {Save the class configuration.};
  \node (ending) [ch4 act decision wide, below=of save] {End class, stage, or term?};
  \node (pastterms) [ch4 act step medium, below=14mm of ending] {Move the class into Past Terms.};
  \node (history) [ch4 act step medium, below=of pastterms] {Preserve class history, membership records, and academic outcomes.};
  \node (end) [edu final, below=of history] {};

  \path [edu arrow] (start.south) -- (open.north);
  \path [edu arrow] (open) -- (role);
  \path [edu arrow] (role.west) -- node[edu label] {Administrator} (adminview.east);
  \path [edu arrow] (role.east) -- node[edu label] {Teacher} (teacherview.west);
  \path [edu arrow] (adminview.south) |- (details.west);
  \path [edu arrow] (teacherview.south) |- (details.east);
  \path [edu arrow] (details) -- (members);
  \path [edu arrow] (members) -- (features);
  \path [edu arrow] (features) -- (save);
  \path [edu arrow] (save) -- (ending);
  \path [edu arrow] (ending) -- node[edu label] {Yes} (pastterms);
  \path [edu arrow] (ending.east) -- ++(10mm,0) node[edu label, pos=0.18] {No} |- (end.east);
  \path [edu arrow] (pastterms) -- (history);
  \path [edu arrow] (history) -- (end);
\end{tikzpicture}%
}

\newcommand{\EduActivityMaterialsWorkflowDiagram}{%
\begin{tikzpicture}[edu activity, x=1mm, y=1mm]
  \node (start) [edu initial] at (0,0) {};
  \node (open) [ch4 rebuilt step wide] at (0,-12) {Class member opens class materials.};
  \node (shown) [ch4 rebuilt step wide] at (0,-25) {Materials are displayed.};
  \node (select) [ch4 rebuilt step wide] at (0,-38) {Select a materials action.};
  \node (fork) [ch4 rebuilt fork, minimum width=126mm] at (0,-51) {};

  \draw[ch4 rebuilt card] (-80,-62) rectangle (-20,-176);
  \node[ch4 rebuilt title] at (-76,-69) {A. Manage material};
  \node (manage) [ch4 rebuilt step narrow] at (-50,-83) {Teacher/Admin manages material.};
  \node (managechoice) [ch4 rebuilt decision] at (-50,-101) {Upload or delete?};
  \node (delete) [ch4 rebuilt step narrow] at (-69,-124) {Delete selected material.};
  \node (details) [ch4 rebuilt step narrow] at (-31,-124) {Enter details and attach file.};
  \node (publish) [ch4 rebuilt step narrow] at (-31,-145) {Publish material.};
  \node (updated) [ch4 rebuilt step] at (-50,-166) {Updated materials are shown.};
  \node (manageend) [edu final] at (-50,-182) {};

  \draw[ch4 rebuilt card] (-17,-62) rectangle (17,-140);
  \node[ch4 rebuilt title] at (-13,-69) {B. View material};
  \node (browse) [ch4 rebuilt step narrow] at (0,-87) {Student browses, searches, or filters materials.};
  \node (preview) [ch4 rebuilt step narrow] at (0,-114) {Student previews, opens, or downloads material.};
  \node (viewend) [edu final] at (0,-135) {};

  \draw[ch4 rebuilt card] (20,-62) rectangle (80,-186);
  \node[ch4 rebuilt title] at (24,-69) {C. AI summary};
  \node (aiopen) [ch4 rebuilt step narrow] at (50,-83) {Choose AI support.};
  \node (aiq) [ch4 rebuilt decision] at (50,-101) {AI summary enabled?};
  \node (blocked) [ch4 rebuilt step narrow] at (31,-124) {Unavailable message is shown.};
  \node (summary) [ch4 rebuilt step narrow] at (69,-124) {Generate or open summary.};
  \node (aiout) [ch4 rebuilt step narrow] at (69,-154) {AI output is displayed.};
  \node (blockedend) [edu final] at (31,-154) {};
  \node (aiend) [edu final] at (69,-181) {};

  \path [edu arrow] (start.south) -- (open.north);
  \path [edu arrow] (open) -- (shown);
  \path [edu arrow] (shown) -- (select);
  \path [edu arrow] (select) -- (fork);
  \path [edu arrow] (fork) -| (manage.north);
  \path [edu arrow] (fork) -- (browse.north);
  \path [edu arrow] (fork) -| (aiopen.north);
  \path [edu arrow] (manage) -- (managechoice);
  \path [edu arrow] (managechoice.west) -| node[edu label, pos=0.25] {Delete} (delete.north);
  \path [edu arrow] (managechoice.east) -| node[edu label, pos=0.25] {Upload} (details.north);
  \path [edu arrow] (details) -- (publish);
  \path [edu arrow] (publish) -- (updated);
  \path [edu arrow] (delete.south) |- (updated.west);
  \path [edu arrow] (updated) -- (manageend);

  \path [edu arrow] (browse) -- (preview);
  \path [edu arrow] (preview) -- (viewend);

  \path [edu arrow] (aiopen) -- (aiq);
  \path [edu arrow] (aiq.west) -| node[edu label, pos=0.25] {No} (blocked.north);
  \path [edu arrow] (blocked) -- (blockedend);
  \path [edu arrow] (aiq.east) -| node[edu label, pos=0.25] {Yes} (summary.north);
  \path [edu arrow] (summary) -- (aiout);
  \path [edu arrow] (aiout) -- (aiend);
\end{tikzpicture}%
}

\newcommand{\EduActivityCommunicationNotificationsDiagram}{%
\begin{tikzpicture}[edu activity, x=1mm, y=1mm]
  \node (start) [edu initial] at (0,0) {};
  \node (open) [ch4 rebuilt step wide] at (0,-12) {User opens the class communication area.};
  \node (select) [ch4 rebuilt step wide] at (0,-26) {Select communication action.};
  \node (fork) [ch4 rebuilt fork, minimum width=118mm] at (0,-39) {};
  \node (lowerfork) [ch4 rebuilt fork, minimum width=58mm] at (0,-164) {};

  \draw[ch4 rebuilt card] (-74,-50) rectangle (-5,-151);
  \node[ch4 rebuilt title] at (-70,-57) {A. Send message};
  \node (msgtype) [ch4 rebuilt decision] at (-39,-76) {Choose message type.};
  \node (text) [ch4 rebuilt step narrow] at (-58,-101) {Send text message.};
  \node (textout) [ch4 rebuilt step narrow] at (-58,-126) {Conversation is updated.};
  \node (textend) [edu final] at (-58,-146) {};
  \node (media) [ch4 rebuilt step narrow] at (-20,-101) {Send media message.};
  \node (mediaout) [ch4 rebuilt step narrow] at (-20,-126) {Conversation is updated.};
  \node (mediaend) [edu final] at (-20,-146) {};

  \draw[ch4 rebuilt card] (5,-50) rectangle (74,-151);
  \node[ch4 rebuilt title] at (9,-57) {D. Notifications};
  \node (notif) [ch4 rebuilt step] at (39,-73) {Open notifications.};
  \node (notifview) [ch4 rebuilt step] at (39,-96) {View notification or related item.};
  \node (notifmanage) [ch4 rebuilt step] at (39,-119) {Mark as read, mark all, or delete if needed.};
  \node (notifout) [ch4 rebuilt step] at (39,-142) {Updated notifications are shown.};
  \node (notifend) [edu final] at (39,-158) {};

  \draw[ch4 rebuilt card] (-74,-168) rectangle (-5,-223);
  \node[ch4 rebuilt title] at (-70,-175) {B. Post announcement};
  \node (announce) [ch4 rebuilt step] at (-39,-191) {Post announcement if authorized.};
  \node (announceout) [ch4 rebuilt step] at (-39,-214) {Announcement or conversation is updated.};
  \node (announceend) [edu final] at (-39,-230) {};

  \draw[ch4 rebuilt card] (5,-168) rectangle (74,-223);
  \node[ch4 rebuilt title] at (9,-175) {C. Search conversation};
  \node (search) [ch4 rebuilt step] at (39,-191) {Search conversation.};
  \node (searchout) [ch4 rebuilt step] at (39,-214) {Matching messages are displayed.};
  \node (searchend) [edu final] at (39,-230) {};

  \path [edu arrow] (start.south) -- (open.north);
  \path [edu arrow] (open) -- (select);
  \path [edu arrow] (select) -- (fork);
  \path [edu arrow] (fork) -| (msgtype.north);
  \path [edu arrow] (fork) -| (notif.north);
  \path [edu arrow] (fork) -- (lowerfork);
  \path [edu arrow] (lowerfork.west) -| (announce.north);
  \path [edu arrow] (lowerfork.east) -| (search.north);
  \path [edu arrow] (msgtype.west) -| node[edu label, pos=0.25] {Text} (text.north);
  \path [edu arrow] (text) -- (textout);
  \path [edu arrow] (textout) -- (textend);
  \path [edu arrow] (msgtype.east) -| node[edu label, pos=0.25] {Media} (media.north);
  \path [edu arrow] (media) -- (mediaout);
  \path [edu arrow] (mediaout) -- (mediaend);

  \path [edu arrow] (announce) -- (announceout);
  \path [edu arrow] (announceout) -- (announceend);

  \path [edu arrow] (search) -- (searchout);
  \path [edu arrow] (searchout) -- (searchend);

  \path [edu arrow] (notif) -- (notifview);
  \path [edu arrow] (notifview) -- (notifmanage);
  \path [edu arrow] (notifmanage) -- (notifout);
  \path [edu arrow] (notifout) -- (notifend);
\end{tikzpicture}%
}

\newcommand{\EduActivityAssignmentWorkflowDiagram}{%
\begin{tikzpicture}[edu activity, x=1mm, y=1mm]
  \node (start) [edu initial] at (0,0) {};
  \draw[ch4 rebuilt card] (-62,-8) rectangle (62,-70);
  \node[ch4 rebuilt title] at (-58,-15) {Teacher/Admin authoring};
  \node (author) [ch4 rebuilt step wide] at (0,-26) {Teacher or administrator opens Assignments.};
  \node (create) [ch4 rebuilt step wide] at (0,-43) {Create or edit the assignment.};
  \node (configure) [ch4 rebuilt step wide] at (0,-60) {Enter instructions, due date, score, and submission settings.};

  \node (publish) [ch4 rebuilt decision] at (0,-83) {Publish now?};
  \node (draft) [ch4 rebuilt step narrow] at (-49,-83) {Save as draft.};
  \node (draftend) [edu final] at (-49,-105) {};
  \node (visible) [ch4 rebuilt step wide] at (0,-107) {Published assignment is visible to students.};

  \draw[ch4 rebuilt card] (-62,-120) rectangle (62,-185);
  \node[ch4 rebuilt title] at (-58,-127) {Student submission};
  \node (studentopen) [ch4 rebuilt step wide] at (0,-139) {Student opens the published assignment.};
  \node (submit) [ch4 rebuilt step wide] at (0,-158) {Student submits text or file.};
  \node (resubmit) [ch4 rebuilt decision] at (0,-179) {Resubmit before deadline?};
  \node (update) [ch4 rebuilt step narrow] at (49,-179) {Student updates the submission.};

  \draw[ch4 rebuilt card] (-62,-198) rectangle (62,-254);
  \node[ch4 rebuilt title] at (-58,-205) {Grading and feedback};
  \node (review) [ch4 rebuilt step wide] at (0,-216) {Teacher or administrator reviews submissions.};
  \node (grade) [ch4 rebuilt step wide] at (0,-234) {Grade and feedback are recorded.};
  \node (result) [ch4 rebuilt step wide] at (0,-252) {Student views the returned score and feedback.};
  \node (end) [edu final] at (0,-273) {};

  \path [edu arrow] (start.south) -- (author.north);
  \path [edu arrow] (author) -- (create);
  \path [edu arrow] (create) -- (configure);
  \path [edu arrow] (configure) -- (publish);
  \path [edu arrow] (publish.west) -- node[edu label] {No} (draft.east);
  \path [edu arrow] (draft) -- (draftend);
  \path [edu arrow] (publish) -- node[edu label] {Yes} (visible);
  \path [edu arrow] (visible) -- (studentopen);
  \path [edu arrow] (studentopen) -- (submit);
  \path [edu arrow] (submit) -- (resubmit);
  \path [edu arrow] (resubmit.east) -- node[edu label] {Yes} (update.west);
  \path [edu arrow] (resubmit) -- node[edu label] {No} (review);
  \path [edu arrow] (update.south) |- (review.east);
  \path [edu arrow] (review) -- (grade);
  \path [edu arrow] (grade) -- (result);
  \path [edu arrow] (result) -- (end);
\end{tikzpicture}%
}

\newcommand{\EduActivityExamIntegrityWorkflowDiagram}{%
\begin{tikzpicture}[edu activity, x=1mm, y=1mm]
  \node (start) [edu initial] at (0,0) {};
  \draw[ch4 rebuilt card] (-66,-8) rectangle (66,-82);
  \node[ch4 rebuilt title] at (-62,-15) {Teacher/Admin exam setup};
  \node (manage) [ch4 rebuilt step wide] at (0,-26) {Teacher or administrator opens exam management.};
  \node (configure) [ch4 rebuilt step wide] at (0,-45) {Enter exam title, add questions, and configure passcode and timing.};
  \node (publish) [ch4 rebuilt decision] at (0,-67) {Publish exam?};
  \node (draft) [ch4 rebuilt step narrow] at (-51,-67) {Save as draft.};
  \node (draftend) [edu final] at (-51,-90) {};
  \node (available) [ch4 rebuilt step wide] at (0,-91) {Exam is available during the valid window.};

  \draw[ch4 rebuilt card] (-66,-105) rectangle (66,-240);
  \node[ch4 rebuilt title] at (-62,-112) {Student attempt};
  \node (studentopen) [ch4 rebuilt step wide] at (0,-124) {Student opens the active exam.};
  \node (passcode) [ch4 rebuilt decision] at (0,-145) {Passcode required?};
  \node (entercode) [ch4 rebuilt step narrow] at (49,-145) {Enter passcode.};
  \node (attempt) [ch4 rebuilt step wide] at (0,-166) {Start the attempt and follow visible integrity rules.};
  \node (answer) [ch4 rebuilt step wide] at (0,-184) {Answer questions during the active attempt.};
  \node (integrity) [ch4 rebuilt decision] at (0,-205) {Integrity warning occurs?};
  \node (event) [ch4 rebuilt step narrow] at (49,-205) {Warning or event is shown.};
  \node (submit) [ch4 rebuilt step wide] at (0,-230) {Submit the attempt or finish at expiry.};

  \draw[ch4 rebuilt card] (-74,-244) rectangle (74,-389);
  \node[ch4 rebuilt title] at (-62,-251) {Review and release};
  \node (review) [ch4 rebuilt step wide] at (0,-263) {Teacher or administrator reviews attempts and integrity information.};
  \node (grade) [ch4 rebuilt step wide] at (0,-282) {Manual answers are graded and outcome is finalized.};
  \node (postreview) [ch4 rebuilt decision] at (0,-305) {Select post-review action.};
  \node (releaseaction) [ch4 rebuilt step narrow] at (-52,-332) {Release exam results.};
  \node (voidaction) [ch4 rebuilt step narrow] at (0,-332) {Void selected attempt.};
  \node (retakeaction) [ch4 rebuilt step narrow] at (52,-332) {Grant retake.};
  \node (studentresult) [ch4 rebuilt step narrow] at (-52,-358) {Student views released result and answer review.};
  \node (voidmessage) [ch4 rebuilt step narrow] at (0,-358) {Void status is shown for the attempt.};
  \node (retakeavailable) [ch4 rebuilt step narrow] at (52,-358) {Student can open a new attempt.};
  \node (releaseend) [edu final] at (-52,-381) {};
  \node (voidend) [edu final] at (0,-381) {};

  \path [edu arrow] (start.south) -- (manage.north);
  \path [edu arrow] (manage) -- (configure);
  \path [edu arrow] (configure) -- (publish);
  \path [edu arrow] (publish.west) -- node[edu label] {No} (draft.east);
  \path [edu arrow] (draft) -- (draftend);
  \path [edu arrow] (publish) -- node[edu label] {Yes} (available);
  \path [edu arrow] (available) -- (studentopen);
  \path [edu arrow] (studentopen) -- (passcode);
  \path [edu arrow] (passcode.east) -- node[edu label] {Yes} (entercode.west);
  \path [edu arrow] (passcode) -- node[edu label] {No} (attempt);
  \path [edu arrow] (entercode.south) |- (attempt.east);
  \path [edu arrow] (attempt) -- (answer);
  \path [edu arrow] (answer) -- (integrity);
  \path [edu arrow] (integrity) -- node[edu label] {No} (submit);
  \path [edu arrow] (integrity.east) -- node[edu label] {Yes} (event.west);
  \path [edu arrow] (event.south) |- (submit.east);
  \path [edu arrow] (submit) -- (review);
  \path [edu arrow] (review) -- (grade);
  \path [edu arrow] (grade) -- (postreview);
  \path [edu arrow] (postreview.west) -| node[edu label, pos=0.25] {Release} (releaseaction.north);
  \path [edu arrow] (postreview) -- node[edu label] {Void} (voidaction);
  \path [edu arrow] (postreview.east) -| node[edu label, pos=0.25] {Retake} (retakeaction.north);
  \path [edu arrow] (releaseaction) -- (studentresult);
  \path [edu arrow] (studentresult) -- (releaseend);
  \path [edu arrow] (voidaction) -- (voidmessage);
  \path [edu arrow] (voidmessage) -- (voidend);
  \path [edu arrow] (retakeaction) -- (retakeavailable);
  \path [edu arrow] (retakeavailable.east) -- ++(18mm,0) |- node[edu label, pos=0.78] {New attempt} (studentopen.east);
\end{tikzpicture}%
}

\newcommand{\EduActivityLiveSessionWorkflowDiagram}{%
\begin{tikzpicture}[edu activity, x=1mm, y=1mm]
  \node (start) [edu initial] at (0,0) {};
  \node (open) [ch4 rebuilt step wide] at (0,-13) {Teacher opens the Sessions page.};
  \node (startsession) [ch4 rebuilt step wide] at (0,-29) {Teacher starts the live session for the selected class.};
  \node (available) [ch4 rebuilt step wide] at (0,-45) {Live session is available to enrolled students.};
  \node (fork) [ch4 rebuilt fork, minimum width=114mm] at (0,-59) {};

  \draw[ch4 rebuilt card] (-74,-70) rectangle (-5,-173);
  \node[ch4 rebuilt title] at (-70,-77) {Teacher controls};
  \node (teacherjoin) [ch4 rebuilt step] at (-39,-93) {Teacher joins and manages microphone and camera.};
  \node (share) [ch4 rebuilt step] at (-39,-119) {Teacher uses screen sharing when needed.};
  \node (teachertools) [ch4 rebuilt step] at (-39,-145) {Teacher uses participants list, session chat, and whiteboard tools.};
  \node (endingsession) [ch4 rebuilt step] at (-39,-166) {Teacher ends the live session.};
  \node (teacherend) [edu final] at (-39,-186) {};

  \draw[ch4 rebuilt card] (5,-70) rectangle (74,-173);
  \node[ch4 rebuilt title] at (9,-77) {Student participation};
  \node (studentjoin) [ch4 rebuilt step] at (39,-93) {Enrolled student joins the live session.};
  \node (studentchat) [ch4 rebuilt step] at (39,-119) {Student views participants and uses session chat.};
  \node (studentwhiteboard) [ch4 rebuilt step] at (39,-145) {Student views the shared whiteboard.};
  \node (studentleave) [ch4 rebuilt step] at (39,-166) {Student leaves or the session ends.};
  \node (studentend) [edu final] at (39,-186) {};

  \path [edu arrow] (start.south) -- (open.north);
  \path [edu arrow] (open) -- (startsession);
  \path [edu arrow] (startsession) -- (available);
  \path [edu arrow] (available) -- (fork);
  \path [edu arrow] (fork.west) -| node[edu label, pos=0.25] {Teacher} (teacherjoin.north);
  \path [edu arrow] (fork.east) -| node[edu label, pos=0.25] {Student} (studentjoin.north);
  \path [edu arrow] (teacherjoin) -- (share);
  \path [edu arrow] (share) -- (teachertools);
  \path [edu arrow] (teachertools) -- (endingsession);
  \path [edu arrow] (endingsession) -- (teacherend);
  \path [edu arrow] (studentjoin) -- (studentchat);
  \path [edu arrow] (studentchat) -- (studentwhiteboard);
  \path [edu arrow] (studentwhiteboard) -- (studentleave);
  \path [edu arrow] (studentleave) -- (studentend);
\end{tikzpicture}%
}

\newcommand{\EduActivityResultsWorkflowDiagram}{%
\begin{tikzpicture}[edu activity, x=1mm, y=1mm]
  \node (start) [edu initial] at (0,0) {};
  \node (open) [ch4 rebuilt step wide] at (0,-12) {User opens the Results area.};
  \node (select) [ch4 rebuilt step wide] at (0,-28) {Select result or score view.};
  \node (fork) [ch4 rebuilt fork, minimum width=94mm] at (0,-42) {};

  \draw[ch4 rebuilt card] (-74,-54) rectangle (-4,-196);
  \node[ch4 rebuilt title] at (-70,-61) {Student branch};
  \node (released) [ch4 rebuilt decision] at (-39,-78) {Results released or visible?};
  \node (hidden) [ch4 rebuilt step narrow] at (-58,-107) {Unavailable results remain hidden or marked not released.};
  \node (hiddenend) [edu final] at (-58,-132) {};
  \node (ownresults) [ch4 rebuilt step narrow] at (-20,-107) {View own scores, grades, and feedback.};
  \node (examreview) [ch4 rebuilt step narrow] at (-20,-134) {View released exam review where allowed.};
  \node (sharedq) [ch4 rebuilt decision] at (-20,-160) {Classmate scores enabled?};
  \node (sharedscores) [ch4 rebuilt step narrow] at (-55,-182) {View shared class scores/results.};
  \node (studentend) [edu final] at (-20,-205) {};

  \draw[ch4 rebuilt card] (4,-54) rectangle (74,-158);
  \node[ch4 rebuilt title] at (8,-61) {Teacher/Admin branch};
  \node (staffopen) [ch4 rebuilt step] at (39,-83) {Open class, assignment, or exam results.};
  \node (staffreview) [ch4 rebuilt step] at (39,-112) {Review results.};
  \node (staffend) [edu final] at (39,-144) {};

  \path [edu arrow] (start.south) -- (open.north);
  \path [edu arrow] (open) -- (select);
  \path [edu arrow] (select) -- (fork);
  \path [edu arrow] (fork) -| node[edu label, pos=0.25] {Student} (released.north);
  \path [edu arrow] (fork) -| node[edu label, pos=0.25] {Teacher/Admin} (staffopen.north);
  \path [edu arrow] (released.west) -| node[edu label, pos=0.25] {No} (hidden.north);
  \path [edu arrow] (hidden) -- (hiddenend);
  \path [edu arrow] (released.east) -| node[edu label, pos=0.25] {Yes} (ownresults.north);
  \path [edu arrow] (ownresults) -- (examreview);
  \path [edu arrow] (examreview) -- (sharedq);
  \path [edu arrow] (sharedq.west) -| node[edu label, pos=0.25] {Yes} (sharedscores.north);
  \path [edu arrow] (sharedscores.south) |- (studentend.west);
  \path [edu arrow] (sharedq) -- node[edu label] {No} (studentend);

  \path [edu arrow] (staffopen) -- (staffreview);
  \path [edu arrow] (staffreview) -- (staffend);
\end{tikzpicture}%
}

\newcommand{\EduActivityDashboardWorkflowDiagram}{%
\begin{tikzpicture}[edu activity, x=1mm, y=1mm]
  \node (start) [edu initial] at (0,0) {};
  \node (open) [ch4 rebuilt step wide] at (0,-12) {User opens the dashboard.};
  \node (role) [ch4 rebuilt step wide] at (0,-28) {Active role and organization context are shown.};
  \node (fork) [ch4 rebuilt fork, minimum width=118mm] at (0,-43) {};

  \draw[ch4 rebuilt card] (-74,-55) rectangle (-26,-151);
  \node[ch4 rebuilt title] at (-70,-62) {Student dashboard};
  \node (studentdash) [ch4 rebuilt step narrow] at (-50,-80) {Student dashboard overview is displayed.};
  \node (studentitems) [ch4 rebuilt step narrow] at (-50,-110) {Review active classes, tasks, notifications, and learning links.};
  \node (studentend) [edu final] at (-50,-142) {};

  \draw[ch4 rebuilt card] (-22,-55) rectangle (22,-151);
  \node[ch4 rebuilt title] at (-18,-62) {Teacher dashboard};
  \node (teacherdash) [ch4 rebuilt step narrow] at (0,-80) {Teacher dashboard overview is displayed.};
  \node (teacheritems) [ch4 rebuilt step narrow] at (0,-110) {Review classes, grading workload, sessions, and class activity.};
  \node (teacherend) [edu final] at (0,-142) {};

  \draw[ch4 rebuilt card] (26,-55) rectangle (74,-151);
  \node[ch4 rebuilt title] at (30,-62) {Admin dashboard};
  \node (admindash) [ch4 rebuilt step narrow] at (50,-80) {Admin dashboard overview is displayed.};
  \node (adminitems) [ch4 rebuilt step narrow] at (50,-110) {Review organization statistics, members, classes, and settings summary.};
  \node (adminend) [edu final] at (50,-142) {};

  \path [edu arrow] (start.south) -- (open.north);
  \path [edu arrow] (open) -- (role);
  \path [edu arrow] (role) -- (fork);
  \path [edu arrow] (fork.west) -| node[edu label, pos=0.25] {Student} (studentdash.north);
  \path [edu arrow] (fork) -- node[edu label] {Teacher} (teacherdash);
  \path [edu arrow] (fork.east) -| node[edu label, pos=0.25] {Admin} (admindash.north);
  \path [edu arrow] (studentdash) -- (studentitems);
  \path [edu arrow] (studentitems) -- (studentend);
  \path [edu arrow] (teacherdash) -- (teacheritems);
  \path [edu arrow] (teacheritems) -- (teacherend);
  \path [edu arrow] (admindash) -- (adminitems);
  \path [edu arrow] (adminitems) -- (adminend);
\end{tikzpicture}%
}

\newcommand{\EduActivityAiLearningSupportWorkflowDiagram}{%
\resizebox{!}{0.74\textheight}{%
\begin{tikzpicture}[edu activity, x=1mm, y=1mm]
  \draw[ch4 rebuilt card] (-60,5) rectangle (60,-58);
  \node[ch4 rebuilt title] at (-56,0) {AI availability};
  \node (start) [edu initial] at (0,-4) {};
  \node (select) [ch4 rebuilt step wide] at (0,-16) {User selects an AI-related action.};
  \node (enabled) [ch4 rebuilt decision] at (0,-32) {AI enabled for current organization and class?};
  \node (blocked) [ch4 rebuilt step narrow] at (40,-32) {AI unavailable message is shown.};
  \node (blockedend) [edu final] at (40,-51) {};
  \node (category) [ch4 rebuilt step wide] at (0,-52) {Choose AI support category.};
  \node (fork) [ch4 rebuilt fork, minimum width=118mm] at (0,-67) {};

  \draw[ch4 rebuilt card] (-74,-76) rectangle (-6,-137);
  \node[ch4 rebuilt title] at (-70,-83) {A. Class AI Agent};
  \node (ask) [ch4 rebuilt step] at (-40,-98) {Ask the Class AI Agent.};
  \node (askout) [ch4 rebuilt step] at (-40,-120) {AI response is displayed using the current class context.};
  \node (askend) [edu final] at (-40,-134) {};

  \draw[ch4 rebuilt card] (6,-76) rectangle (74,-137);
  \node[ch4 rebuilt title] at (10,-83) {B. Study support};
  \node (summary) [ch4 rebuilt step] at (40,-98) {Generate material summary.};
  \node (studyout) [ch4 rebuilt step] at (40,-120) {AI study support is displayed using the current class context.};
  \node (studyend) [edu final] at (40,-134) {};

  \draw[ch4 rebuilt card] (-74,-151) rectangle (74,-232);
  \node[ch4 rebuilt title] at (-70,-158) {C. Drafting and feedback support};
  \node (request) [ch4 rebuilt step wide] at (0,-173) {Request assignment support or drafting assistance.};
  \node (task) [ch4 rebuilt decision] at (0,-192) {Selected support type?};
  \node (support) [ch4 rebuilt step narrow] at (-48,-215) {Request assignment support.};
  \node (draft) [ch4 rebuilt step narrow] at (0,-215) {Generate assignment or exam draft if enabled.};
  \node (feedback) [ch4 rebuilt step narrow] at (48,-215) {Generate draft feedback if supported.};
  \node (draftout) [ch4 rebuilt step wide] at (0,-232) {Draft AI output is displayed for user review.};
  \node (draftend) [edu final] at (0,-247) {};

  \path [edu arrow] (start.south) -- (select.north);
  \path [edu arrow] (select) -- (enabled);
  \path [edu arrow] (enabled.east) -- node[edu label] {No} (blocked.west);
  \path [edu arrow] (blocked) -- (blockedend);
  \path [edu arrow] (enabled) -- node[edu label] {Yes} (category);

  \path [edu arrow] (category) -- (fork);
  \path [edu arrow] (fork) -| (ask.north);
  \path [edu arrow] (fork) -| (summary.north);
  \path [edu arrow] (fork) -- (request.north);
  \path [edu arrow] (ask) -- (askout);
  \path [edu arrow] (askout) -- (askend);

  \path [edu arrow] (summary) -- (studyout);
  \path [edu arrow] (studyout) -- (studyend);

  \path [edu arrow] (request) -- (task);
  \path [edu arrow] (task.west) -| node[edu label, pos=0.25] {Support} (support.north);
  \path [edu arrow] (task) -- node[edu label] {Assignment/exam} (draft);
  \path [edu arrow] (task.east) -| node[edu label, pos=0.25] {Feedback} (feedback.north);
  \path [edu arrow] (support.south) |- (draftout.west);
  \path [edu arrow] (draft) -- (draftout);
  \path [edu arrow] (feedback.south) |- (draftout.east);
  \path [edu arrow] (draftout) -- (draftend);
\end{tikzpicture}%
}%
}

\newcommand{\EduActivityIdeWorkflowDiagram}{%
\begin{tikzpicture}[edu activity, node distance=8.5mm and 28mm]
  \node (start) [edu initial] {};
  \node (open) [ch4 act step] {User opens the Integrated Coding Lab / IDE workspace.};
  \node (environment) [ch4 act step medium, below=of open] {Choose the available coding environment or extension for the class.};
  \node (workspace) [ch4 act step medium, below=of environment] {System loads the workspace and visible project files.};
  \node (files) [ch4 act step medium, below=of workspace] {Browse, create, rename, delete, or open files as permitted.};
  \node (edit) [ch4 act step medium, below=of files] {Edit code or markdown content and save the current work.};
  \node (run) [ch4 act decision wide, below=of edit] {Run or preview the current work?};
  \node (execute) [ch4 act step medium, below=14mm of run] {Run the code and open the available preview.};
  \node (inspect) [ch4 act step medium, below=of execute] {Review the preview, console output, or problems panel.};
  \node (extras) [ch4 act step medium, below=of inspect] {Use the IDE terminal where supported or reset the workspace when needed.};
  \node (end) [edu final, below=of extras] {};

  \path [edu arrow] (start.south) -- (open.north);
  \path [edu arrow] (open) -- (environment);
  \path [edu arrow] (environment) -- (workspace);
  \path [edu arrow] (workspace) -- (files);
  \path [edu arrow] (files) -- (edit);
  \path [edu arrow] (edit) -- (run);
  \path [edu arrow] (run) -- node[edu label] {Yes} (execute);
  \path [edu arrow] (run.east) -- ++(10mm,0) node[edu label, pos=0.18] {No} |- (end.east);
  \path [edu arrow] (execute) -- (inspect);
  \path [edu arrow] (inspect) -- (extras);
  \path [edu arrow] (extras) -- (end);
\end{tikzpicture}%
}

\newcommand{\EduActivityProfileHelpWorkflowDiagram}{%
\begin{tikzpicture}[edu activity, x=1mm, y=1mm]
  \node (start) [edu initial] at (0,0) {};
  \node (open) [ch4 rebuilt step wide] at (0,-12) {User opens account controls or the help area.};
  \node (route) [ch4 rebuilt decision] at (0,-31) {Select support action.};
  \node (fork) [ch4 rebuilt fork, minimum width=104mm] at (0,-47) {};

  \draw[ch4 rebuilt card] (-74,-58) rectangle (-5,-190);
  \node[ch4 rebuilt title] at (-70,-65) {A. Profile and workspace};
  \node (profile) [ch4 rebuilt step] at (-39,-83) {View profile information.};
  \node (settings) [ch4 rebuilt step] at (-39,-110) {Update implemented profile or display settings.};
  \node (switch) [ch4 rebuilt step] at (-39,-137) {Switch active organization or role if available.};
  \node (profileout) [ch4 rebuilt step] at (-39,-160) {Workspace context is refreshed.};
  \node (profileend) [edu final] at (-39,-182) {};

  \draw[ch4 rebuilt card] (5,-58) rectangle (74,-166);
  \node[ch4 rebuilt title] at (9,-65) {B. Help center};
  \node (help) [ch4 rebuilt step] at (39,-83) {Open the help center.};
  \node (search) [ch4 rebuilt step] at (39,-110) {Search or browse help articles.};
  \node (article) [ch4 rebuilt step] at (39,-137) {Open the selected article.};
  \node (helpend) [edu final] at (39,-159) {};

  \path [edu arrow] (start.south) -- (open.north);
  \path [edu arrow] (open) -- (route);
  \path [edu arrow] (route) -- (fork);
  \path [edu arrow] (fork.west) -| node[edu label, pos=0.25] {Profile} (profile.north);
  \path [edu arrow] (fork.east) -| node[edu label, pos=0.25] {Help} (help.north);
  \path [edu arrow] (profile) -- (settings);
  \path [edu arrow] (settings) -- (switch);
  \path [edu arrow] (switch) -- (profileout);
  \path [edu arrow] (profileout) -- (profileend);
  \path [edu arrow] (help) -- (search);
  \path [edu arrow] (search) -- (article);
  \path [edu arrow] (article) -- (helpend);
\end{tikzpicture}%
}


\newcommand{\EduUseCasePdfFigure}[3]{%
  \begin{figure}[htbp]
    \centering
    \includegraphics[page=1,width=\textwidth,height=0.82\textheight,keepaspectratio]{#1}
    \caption{#2}
    \label{#3}
  \end{figure}
}

\section{Introduction}\label{sec:system_analysis_introduction}

This chapter translates the strategic direction of EduVerse into its principal
analytical elements. It identifies the system actors, functional and
non-functional requirements, access conditions, and interaction models that
define expected platform behavior.

The chapter then presents the use-case and activity models used to represent
user interaction and major workflow sequences. Together, these elements
provide the analytical basis for the system design and implementation described
in the following chapters.

\section{System Overview}\label{sec:system_overview}

EduVerse is a multi-organization educational platform in which the class acts
as the primary operational workspace. The system supports three principal
human roles: Organization Administrator, Teacher, and Student, while the same
account may belong to multiple organizations or hold more than one role.

System behavior depends on the active organization, active role, class
membership, and applicable feature configuration. These conditions determine
which resources, workflows, and management actions are available to the
current user.

The web application provides the primary operational environment, while the
mobile companion extends selected high-frequency academic interactions. This
chapter analyzes the system at the level of actors, requirements, permissions,
and workflows rather than implementation architecture.

\section{System Actors}\label{sec:system_actors}

The principal human actors of EduVerse are the Organization Administrator,
Teacher, and Student. These roles represent different levels of responsibility
within the platform rather than permanently separate account types. A user may
belong to more than one organization or hold more than one role, with access
resolved according to the current workspace context.

\begin{table}[htbp]
  \centering
  \small
  \renewcommand{\arraystretch}{1.15}
  \setlength{\tabcolsep}{5pt}
  \caption{Primary EduVerse System Actors}
  \label{tab:eduverse_system_actors}
  \begin{tabular}{|p{0.23\textwidth}|p{0.57\textwidth}|}
    \hline
    \textbf{Actor} & \textbf{Analytical Role in the System} \\
    \hline
    Organization Administrator & Manages organization-level governance, membership, role assignment, institutional settings, class oversight, and other administrative capabilities. \\
    \hline
    Teacher & Manages instructional and class-level academic workflows within the permissions granted by the organization and class context. \\
    \hline
    Student & Participates in permitted learning activities, accesses class resources, completes assessments, joins communication and live-learning workflows, and reviews released academic outcomes. \\
    \hline
  \end{tabular}
\end{table}

The diagrams in this chapter focus on the three principal human actors.
Automatic responsibilities such as access enforcement, notification generation,
integrity-event recording, and workflow-state changes are represented through
the surrounding analysis and workflow behavior where relevant.

\section{Functional Requirements}\label{sec:functional_requirements}

The functional requirements define the main behaviors that EduVerse must
provide within the implemented project scope. They are grouped by functional
area to maintain traceability with the use-case model and system workflows.

\subsection*{Authentication and Access}

\begin{itemize}
  \item \textbf{FR-01.} The system shall allow users to sign up, sign in,
    verify email addresses, request password reset, and log out.
  \item \textbf{FR-02.} The system shall require authenticated access before
    opening organization-scoped or class-scoped workspaces.
  \item \textbf{FR-03.} The system shall allow users with more than one
    eligible workspace context to select the Active Organization and Active
    Role.
  \item \textbf{FR-04.} The system shall open the role-appropriate workspace
    or dashboard after successful authentication and context resolution.
\end{itemize}

\subsection*{Organization and Role Management}

\begin{itemize}
  \item \textbf{FR-05.} The system shall allow an Organization Administrator
    to create an organization from a selected template and configure its
    initial institutional settings.
  \item \textbf{FR-06.} The system shall allow Organization Administrators to
    invite or reinvite members, revoke outstanding invitations, manage
    organization membership, and grant or remove permitted organization roles.
  \item \textbf{FR-07.} The system shall support invitation-based access and,
    when enabled, Public Join Link access with target-role selection and
    direct-join or approval-based entry modes.
  \item \textbf{FR-08.} The system shall allow Organization Administrators to
    manage organization-level access governance, including join-request
    decisions, public-access settings, and teacher class permissions.
\end{itemize}

\subsection*{Class Management and Class History}

\begin{itemize}
  \item \textbf{FR-09.} The system shall allow authorized users to create and
    edit classes with details such as title, code, term, stage, description,
    purpose, and related settings.
  \item \textbf{FR-10.} The system shall allow class membership management,
    including assigning or changing the responsible Teacher and managing
    student enrollment according to permissions.
  \item \textbf{FR-11.} The system shall allow authorized users to configure
    class tools, extension availability, visibility rules, Results visibility,
    and class purpose within the controls exposed for the class.
  \item \textbf{FR-12.} The system shall allow administrators or permitted
    teachers to end a class, stage, or term and move the class into Past
    Terms while preserving class history and authorized previous-term student
    linkage.
\end{itemize}

\subsection*{Course Content and Materials}

\begin{itemize}
  \item \textbf{FR-13.} The system shall allow authorized users to upload
    materials together with descriptive details and the selected file.
  \item \textbf{FR-14.} The system shall allow class members to view, search,
    filter, preview, open, and download available materials according to
    access permissions.
  \item \textbf{FR-15.} The system shall allow authorized users to delete
    materials from the applicable class workspace.
\end{itemize}

\subsection*{Class Communication and Notifications}

\begin{itemize}
  \item \textbf{FR-16.} The system shall allow class members to send text or
    media messages within the class communication space.
  \item \textbf{FR-17.} The system shall allow authorized users to post
    announcements, remove announcement entries from the announcement area, and
    search the class conversation.
  \item \textbf{FR-18.} The system shall generate notifications from relevant
    system events and shall allow users to view, open, mark, mark all, and
    delete notification entries.
\end{itemize}

\subsection*{Assignment Management}

\begin{itemize}
  \item \textbf{FR-19.} The system shall allow authorized users to create,
    edit, delete, save as draft, and publish assignments.
  \item \textbf{FR-20.} The system shall allow assignment configuration for
    due date, score range, submission modes, and optional attached files.
  \item \textbf{FR-21.} The system shall allow students to submit text, file,
    or combined work according to assignment settings and to resubmit when the
    assignment policy permits it.
  \item \textbf{FR-22.} The system shall allow Teachers to review submissions,
    record grades, provide feedback, and use AI-assisted feedback tools when
    enabled. The system shall allow Organization Administrators to review
    submissions and grading outcomes for oversight.
\end{itemize}

\subsection*{Exam Management and Integrity}

\begin{itemize}
  \item \textbf{FR-23.} The system shall allow authorized users to create,
    edit, delete, draft, and publish exams, including question management and
    passcode configuration when required.
  \item \textbf{FR-24.} The system shall allow eligible students to start an
    exam attempt, enter a passcode when required, answer questions, and submit
    the attempt within the valid exam window.
  \item \textbf{FR-25.} The system shall apply configured integrity controls,
    including fullscreen requirements and integrity-event recording when the
    required exam state is interrupted.
  \item \textbf{FR-26.} The system shall allow teachers and administrators to
    monitor attempts, review integrity events, grade manual answers, void
    attempts, and grant retakes where permitted.
  \item \textbf{FR-27.} The system shall allow release of exam results and,
    after release, shall allow students to view result details and answer
    review where the class settings permit it.
\end{itemize}

\subsection*{Live Sessions and Whiteboard}

\begin{itemize}
  \item \textbf{FR-28.} The system shall allow authorized class managers to
    start, manage, and end live sessions.
  \item \textbf{FR-29.} The system shall allow eligible participants to join
    a live session and use session tools such as microphone, camera, screen
    sharing, participants view, and session chat.
  \item \textbf{FR-30.} The system shall allow whiteboard interaction or
    whiteboard viewing according to the role and permissions of the current
    participant.
\end{itemize}

\subsection*{Results and Dashboards}

\begin{itemize}
  \item \textbf{FR-31.} The system shall provide role-specific dashboards for
    Students, Teachers, and Organization Administrators.
  \item \textbf{FR-32.} The system shall allow students to view assignment
    grades, feedback, released exam results, and related result details within
    permitted classes.
  \item \textbf{FR-33.} The system shall allow teachers and administrators to
    review class results, assignment results, and exam results subject to the
    configured visibility rules.
\end{itemize}

\subsection*{AI Learning Support}

\begin{itemize}
  \item \textbf{FR-34.} The system shall provide a class AI Agent for
    contextual academic support when the AI feature is enabled for the current
    organization and class context.
  \item \textbf{FR-35.} The system shall allow AI-assisted study support such
    as material summaries and assignment-related guidance without replacing the
    original material or teacher judgment.
  \item \textbf{FR-36.} The system shall allow authorized Teachers and
    Organization Administrators to use AI-assisted drafting for assignments,
    examinations, and assignment feedback when AI support is available in the
    current context.
\end{itemize}

\subsection*{Integrated Coding Lab / IDE Extension}

\begin{itemize}
  \item \textbf{FR-37.} The system shall allow users to open the Integrated
    Coding Lab / IDE Extension when it is enabled for the current class.
  \item \textbf{FR-38.} The system shall allow users to choose an available
    coding environment and manage the visible project files in the workspace.
  \item \textbf{FR-39.} The system shall allow users to open, edit, and save
    code, and to run or preview supported code according to the selected coding
    environment.
  \item \textbf{FR-40.} The system shall allow users to review console or
    problems output, reset the workspace when needed, and use the IDE terminal
    where that capability is supported.
\end{itemize}

\subsection*{Feature Configuration and Extensibility}

\begin{itemize}
  \item \textbf{FR-41.} The system shall allow supported platform features to
    be enabled or disabled at the organization level and at the class level for
    features that provide class-level control.
  \item \textbf{FR-42.} The system shall hide, block, or clearly mark
    unavailable features according to the active organization and class
    configuration rather than exposing unusable controls.
  \item \textbf{FR-43.} The system shall allow users to open enabled extensions
    and feature modules from the applicable organization or class context.
\end{itemize}

\subsection*{Profile and Help}

\begin{itemize}
  \item \textbf{FR-44.} The system shall allow users to view profile
    information, change theme preferences, and manage password-related account
    settings.
  \item \textbf{FR-45.} The system shall allow authenticated users to return to
    workspace selection and change the Active Organization or Active Role
    without signing out.
  \item \textbf{FR-46.} The system shall provide searchable help content for
    account, class, administration, AI, coding, assignment, exam, results,
    navigation, and troubleshooting topics.
\end{itemize}

\section{Non-Functional Requirements}\label{sec:non_functional_requirements}

The non-functional requirements define the principal quality, security,
operational, and maintainability constraints expected of EduVerse within the
evaluated project scope.

\begin{itemize}
  \item \textbf{NFR-01 Usability.} The interface shall maintain consistent
    navigation and clear context indication across the evaluated organization,
    role, class, and feature workflows.
  \item \textbf{NFR-02 Security and Access Control.} The system shall enforce
    authenticated, organization-aware, role-aware, and class-aware access
    controls for protected resources. Public or externally reachable
    educational resources shall remain subject to authenticated or
    organization-defined access rules where required.
  \item \textbf{NFR-03 Reliability.} The system shall preserve consistent
    workflow state and stored records during the evaluated membership,
    materials, submission, examination, notification, live-session, and Results
    workflows.
  \item \textbf{NFR-04 Performance.} Within the evaluated project context,
    interactive pages and tested academic workflows shall respond without
    delays that prevent normal task completion.
  \item \textbf{NFR-05 Maintainability.} The system shall keep supported
    features and shared services structurally separated so that correcting or
    extending one workflow does not require unrelated workflows to be
    rewritten.
  \item \textbf{NFR-06 Scalability.} The data and access model shall support
    multiple organizations, classes, memberships, and academic records without
    requiring a separate application instance for each organization.
  \item \textbf{NFR-07 Data Privacy.} Personal and academic data shall be
    exposed only to authorized users and contexts, with particular protection
    for profiles, submissions, integrity records, grades, and
    organization-scoped data.
  \item \textbf{NFR-08 Cross-Client Consistency.} The system shall apply
    consistent authentication, role-resolution, and core access rules across
    supported clients that use the shared platform services.
\end{itemize}

\section{Use-Case Model}\label{sec:use_case_model}

The use-case model represents the principal interactions between EduVerse users
and the services available within the active organization and class context.
The Project Overview provides a high-level system view, while the remaining
diagrams focus on specific functional areas.

The diagrams emphasize human interaction with the platform. Automatic system
behavior is represented through the surrounding analysis and related workflow
descriptions rather than modeled repeatedly as a separate system actor.

\subsection*{Project Overview}

\EduUseCasePdfFigure
  {figures/chapter4/usecases/project_overview.pdf}
  {Project Overview Use-Case Diagram}
  {fig:ch4_usecase_project_overview}

The Project Overview use-case diagram provides a high-level representation of
the relationship between the three principal human actors and the main
EduVerse service areas. Detailed interactions are decomposed into the
module-specific diagrams that follow.

\subsection*{Authentication and Access}

\EduUseCasePdfFigure
  {figures/chapter4/usecases/authentication_access.pdf}
  {Authentication and Access Use-Case Diagram}
  {fig:ch4_usecase_authentication_access}

Figure~\ref{fig:ch4_usecase_authentication_access} represents entry into the
EduVerse workspace and the conditions required to resolve the user's active
context. It distinguishes account access from organization- and role-selection
behavior when multiple valid contexts are available.

\subsection*{Organization and Role Management}

\EduUseCasePdfFigure
  {figures/chapter4/usecases/organization_role_management.pdf}
  {Organization and Role Management Use-Case Diagram}
  {fig:ch4_usecase_organization_role_management}

Figure~\ref{fig:ch4_usecase_organization_role_management} represents
organization-level governance and membership control. It emphasizes how
administrative authority, role assignment, access pathways, and institutional
settings are resolved within an organization context.

\subsection*{Class Management}

\EduUseCasePdfFigure
  {figures/chapter4/usecases/class_management.pdf}
  {Class Management Use-Case Diagram}
  {fig:ch4_usecase_class_management}

Figure~\ref{fig:ch4_usecase_class_management} models the Class as the central
operational unit of EduVerse. It emphasizes class lifecycle, membership,
configuration, visibility, and historical continuity under the applicable
permissions.

\subsection*{Assignment Management}

\EduUseCasePdfFigure
  {figures/chapter4/usecases/assignment_management.pdf}
  {Assignment Management Use-Case Diagram}
  {fig:ch4_usecase_assignment_management}

Figure~\ref{fig:ch4_usecase_assignment_management} models the assignment
lifecycle and distinguishes teacher-managed authoring and grading activities
from student-controlled submission behavior. Conditional relations represent
states such as publication, submission mode, and permitted resubmission.

\subsection*{AI Learning Support}

\EduUseCasePdfFigure
  {figures/chapter4/usecases/ai_learning_support.pdf}
  {AI Learning Support Use-Case Diagram}
  {fig:ch4_usecase_ai_learning_support}

Figure~\ref{fig:ch4_usecase_ai_learning_support} represents AI as a supporting
capability within the academic workflow rather than as an independent human
actor. Availability depends on the applicable organization and class
configuration.

\subsection*{Class Communication and Notifications}

\EduUseCasePdfFigure
  {figures/chapter4/usecases/class_communication_notifications.pdf}
  {Class Communication and Notifications Use-Case Diagram}
  {fig:ch4_usecase_class_communication_notifications}

Figure~\ref{fig:ch4_usecase_class_communication_notifications} distinguishes
user-generated class communication from system-generated notification
behavior. The model separates direct interaction from event-driven awareness
while keeping both within the class context.

\subsection*{Course Content and Materials}

\EduUseCasePdfFigure
  {figures/chapter4/usecases/course_content_materials.pdf}
  {Course Content and Materials Use-Case Diagram}
  {fig:ch4_usecase_course_content_materials}

Figure~\ref{fig:ch4_usecase_course_content_materials} represents the controlled
lifecycle of class learning resources and the distinction between authorized
content management and member access.

\subsection*{Exam Management and Integrity}

\EduUseCasePdfFigure
  {figures/chapter4/usecases/exam_management_integrity.pdf}
  {Exam Management and Integrity Use-Case Diagram}
  {fig:ch4_usecase_exam_management_integrity}

Figure~\ref{fig:ch4_usecase_exam_management_integrity} models the controlled
examination lifecycle and the relationship between attempt eligibility,
integrity monitoring, grading, result release, and post-attempt administrative
decisions.

\subsection*{Live Sessions and Whiteboard}

\EduUseCasePdfFigure
  {figures/chapter4/usecases/live_sessions_whiteboard.pdf}
  {Live Sessions and Whiteboard Use-Case Diagram}
  {fig:ch4_usecase_live_sessions_whiteboard}

Figure~\ref{fig:ch4_usecase_live_sessions_whiteboard} represents synchronous
participation within the class context and the collaboration capabilities
available to eligible participants according to their role and permissions.

\subsection*{Integrated Coding Lab / IDE Extension}

\EduUseCasePdfFigure
  {figures/chapter4/usecases/integrated_coding_lab_ide_extension.pdf}
  {Integrated Coding Lab / IDE Extension Use-Case Diagram}
  {fig:ch4_usecase_integrated_coding_lab_ide_extension}

Figure~\ref{fig:ch4_usecase_integrated_coding_lab_ide_extension} represents
access to the integrated coding workspace as a configurable class capability.
The model emphasizes the relationship between workspace interaction and
feature availability.

\subsection*{Profile and Help}

\EduUseCasePdfFigure
  {figures/chapter4/usecases/profile_help.pdf}
  {Profile and Help Use-Case Diagram}
  {fig:ch4_usecase_profile_help}

Figure~\ref{fig:ch4_usecase_profile_help} represents account-oriented and
user-support interactions that apply across the wider platform rather than to
one specific academic workflow.

\subsection*{Results and Dashboards}

\EduUseCasePdfFigure
  {figures/chapter4/usecases/results_dashboards.pdf}
  {Results and Dashboards Use-Case Diagram}
  {fig:ch4_usecase_results_dashboards}

Figure~\ref{fig:ch4_usecase_results_dashboards} represents role-sensitive
access to dashboard information and academic outcomes. The model emphasizes
that result visibility depends on actor responsibility, result release, and
applicable class settings.

\clearpage

\section{Activity Diagrams}\label{sec:activity_diagrams}

The following activity diagrams model the principal EduVerse workflows
represented by the use-case set. They emphasize user-visible actions, decision
points, and relevant automatic system responses while avoiding unnecessary
implementation detail.

\subsection*{AD-01 Authentication and Workspace Access}

\EduActivityFigureReadable{\EduActivityAuthWorkspaceAccessDiagram}
  {Authentication and Workspace Access Activity Diagram}
  {fig:ch4_activity_auth_workspace_access}

Figure~\ref{fig:ch4_activity_auth_workspace_access} emphasizes the main
decision paths involved in account access, verification, recovery,
workspace-context resolution, and exit from the authenticated session.

\subsection*{AD-02 Organization Member Access and Public Join Link}

\EduActivityFigureReadable{\EduActivityOrganizationAccessDiagram}
  {Organization Member Access and Public Join Link Activity Diagram}
  {fig:ch4_activity_organization_access}

Figure~\ref{fig:ch4_activity_organization_access} emphasizes the alternative
membership-entry paths and the administrative decisions required when access
depends on invitation, public joining, or approval.

\subsection*{AD-03 Class Creation, Configuration, and Past Terms}

\EduActivityFigure{\EduActivityClassLifecycleDiagram}
  {Class Creation, Configuration, and Past Terms Activity Diagram}
  {fig:ch4_activity_class_lifecycle}

Figure~\ref{fig:ch4_activity_class_lifecycle} represents the class lifecycle
from creation and configuration through controlled closure and preservation in
Past Terms.

\subsection*{AD-04 Course Content and Materials}

\EduActivityFigureReadable{\EduActivityMaterialsWorkflowDiagram}
  {Course Content and Materials Activity Diagram}
  {fig:ch4_activity_materials}

Figure~\ref{fig:ch4_activity_materials} represents the main decision paths
between authorized material management and member access to available learning
resources.

\subsection*{AD-05 Class Communication and Notifications}

\EduActivityFigureTall{\EduActivityCommunicationNotificationsDiagram}
  {Class Communication and Notifications Activity Diagram}
  {fig:ch4_activity_communication_notifications}

Figure~\ref{fig:ch4_activity_communication_notifications} distinguishes direct
class communication actions from event-driven notification handling within the
same class context.

\subsection*{AD-06 Assignment Workflow}

\EduActivityFigureReadable{\EduActivityAssignmentWorkflowDiagram}
  {Assignment Workflow Activity Diagram}
  {fig:ch4_activity_assignment_workflow}

Figure~\ref{fig:ch4_activity_assignment_workflow} emphasizes the assignment
lifecycle and the decision points surrounding publication, submission,
permitted resubmission, grading, and returned outcomes.

\subsection*{AD-07 Exam Management and Integrity Workflow}

\EduActivityFigureTall{\EduActivityExamIntegrityWorkflowDiagram}
  {Exam Management and Integrity Workflow Activity Diagram}
  {fig:ch4_activity_exam_integrity}

Figure~\ref{fig:ch4_activity_exam_integrity} emphasizes the principal
examination decision paths involving attempt eligibility, integrity handling,
grading, result release, voiding, and retake outcomes.

\subsection*{AD-08 Live Session and Whiteboard Workflow}

\EduActivityFigure{\EduActivityLiveSessionWorkflowDiagram}
  {Live Session and Whiteboard Workflow Activity Diagram}
  {fig:ch4_activity_live_session}

Figure~\ref{fig:ch4_activity_live_session} represents the controlled
synchronous-session lifecycle from session availability through participation
to closure.

\subsection*{AD-09 Results Workflow}

\EduActivityFigureTall{\EduActivityResultsWorkflowDiagram}
  {Results Workflow Activity Diagram}
  {fig:ch4_activity_results}

Figure~\ref{fig:ch4_activity_results} represents the conditions under which
academic results become available to different actors and how visibility rules
affect subsequent review.

\subsection*{AD-10 Dashboard Workflow}

\EduActivityFigureReadable{\EduActivityDashboardWorkflowDiagram}
  {Dashboard Workflow Activity Diagram}
  {fig:ch4_activity_dashboard}

Figure~\ref{fig:ch4_activity_dashboard} represents role-specific dashboard
access and preserves the distinction between summary information and detailed
Results workflows.

\subsection*{AD-11 AI Learning Support Workflow}

\EduActivityFigureReadable{\EduActivityAiLearningSupportWorkflowDiagram}
  {AI Learning Support Workflow Activity Diagram}
  {fig:ch4_activity_ai_learning_support}

Figure~\ref{fig:ch4_activity_ai_learning_support} represents access to
AI-supported learning actions and the configuration conditions that determine
whether those capabilities are available.

\subsection*{AD-12 Integrated Coding Lab / IDE Workflow}

\EduActivityFigure{\EduActivityIdeWorkflowDiagram}
  {Integrated Coding Lab / IDE Workflow Activity Diagram}
  {fig:ch4_activity_ide_workflow}

Figure~\ref{fig:ch4_activity_ide_workflow} represents the main coding-workspace
interaction flow under the condition that the corresponding class capability
is enabled.

\subsection*{AD-13 Profile and Help Workflow}

\EduActivityFigureReadable{\EduActivityProfileHelpWorkflowDiagram}
  {Profile and Help Workflow Activity Diagram}
  {fig:ch4_activity_profile_help}

Figure~\ref{fig:ch4_activity_profile_help} represents the user-support workflow
for account-oriented settings, workspace-context access, and help resources.

\clearpage

\section{Chapter Summary}\label{sec:chapter4_summary}

This chapter defined the analytical structure of EduVerse through its actor
model, functional and non-functional requirements, use-case model, and activity
workflows. Together, these elements establish the expected system behavior,
access conditions, and principal interaction patterns that inform the design
and implementation described in the following chapters.

